\PassOptionsToPackage{sols}{handout}
\documentclass[main]{subfiles}
\setcounterrefbeforedot{chapter}{m-ws6-1}
\setcounterrefafterdot{section}{m-ws6-1}
\addtocounter{section}{-1}
\usepackage{solutions}

\begin{document}

\ifSubfilesClassLoaded{%
  \chaptermark{\chaptersix}
}{}

\section{Integration by Substitution} \label{ws6-1}

\begin{problemonly}
    \begin{framed}

\textbf{Key Ideas}:

\begin{itemize}[topsep=0pt,partopsep=0pt,parsep=8pt,itemsep=0pt]

\item You should be comfortable with integration by substitution as a
  way of ``undoing'' the Chain Rule.

\item You should begin to develop flexibility in figuring out ``what to
  pick for $u$.''  There is a lot of trial and error involved in this,
  so don't feel discouraged if your first guesses don't work out: that's
  normal!  With practice, you will develop stronger intuition about what
  to pick.

\item You should be comfortable using integration by substitution on
  both definite and indefinite integrals, and you should use appropriate
  notation when dealing with the endpoints of definite integrals.

\item You should keep in mind that you can always check whether you've
  correctly found an antiderivative $F$ of a function $f$ by
  differentiating $F$.

\item You should be able to use substitution to evaluate $\displaystyle
  \int \tan x \,dx$.

\end{itemize}
\end{framed}
\end{problemonly}

\begin{framed}
  \textbf{Some basic antiderivatives.} Here are some of the most basic
  antiderivatives you should know.
\begin{enumerate}[topsep=0pt]
      
    \item If $n \neq -1$, $\displaystyle \int u^n\,du = \fitb{1.5in}{\frac{1}{n
        + 1} u^{n + 1} + C}$

    \item $\displaystyle \int \frac{1}{u}\,du = \fitb{1.5in}{\ln |u| + C}$

    \item $\displaystyle \int \sin u\,du = \fitb{1.5in}{-\cos u + C}$

    \item $\displaystyle \int \cos u\,du = \fitb{1.5in}{\sin u + C}$

    \item $\displaystyle \int e^u\,du = \fitb{1.5in}{e^u + C}$

    \item $\displaystyle \int \frac{1}{1 + u^2}\,du = \fitb{1.5in}{\arctan u +
      C}$
\end{enumerate}
  
\end{framed}

\begin{enumerate}
\item 
  \textbf{Warm-up.}  Can you evaluate the following? Which one can you
  \ul{not} evaluate?
  \probonly{}
    \begin{enumerate}

    \item
      \note{This is basic practice, as well something we might look back
        at in {{{{\#3}}(e)}}.}

      \begin{problem}[\vfill]
        $\displaystyle \int_1^4 \frac{1}{3} \sqrt{u}\,du$
      \end{problem}

      \begin{solution}
        An antiderivative of $\sqrt{u} = u^{1/2}$ is $\frac{2}{3}
        u^{3/2}$, so
        \begin{align*}
          \int_1^4 \frac{1}{3} \sqrt{u}\,du &= \left. \frac{1}{3} \cdot
            \frac{2}{3} u^{3/2} \right|_1^4 \\
          &= \frac{2}{9} \cdot 4^{3/2} - \frac{2}{9} \cdot 1^{3/2} \\
          &= \frac{2}{9} \cdot 8 - \frac{2}{9} \cdot 1 \\
          &= \boxed{\frac{14}{9}}
        \end{align*}
      \end{solution}

    \item

      \note{I want them to see here that, although we need to take the
        Chain Rule into account when we're thinking of an antiderivative
        of $\cos(3x)$, it's pretty easy to do that.}

      \begin{problem}[\vfill]
        $\displaystyle \int \cos (3x)\,dx$
      \end{problem}

      \begin{solution}
        You might guess that $\sin(3x)$ is involved in an antiderivative
        of $\cos(3x)$.  The derivative of $\sin(3x)$ is $3 \cos(3x)$, so
        the derivative of $\frac{1}{3} \sin(3x)$ is $\cos(3x)$.
        Therefore, $\displaystyle \int \cos(3x)\,dx = \boxed{\frac{1}{3}
          \sin(3x) + C}$. 
      \end{solution}

    \item
      \note{I'll let them know that this integral actually can't be done.} 

      \begin{problem}[\vfill]
        $\displaystyle \int \cos(x^2)\,dx$
      \end{problem}

      \begin{solution}
        There is no way to write down this indefinite integral!  It turns
        out that $\cos(x^2)$ is an example of a function whose
        antiderivatives cannot be expressed as a finite sum of familiar
        functions (like exponentials, logs, trig functions, etc.).  This
        doesn't mean that $\cos(x^2)$ doesn't have antiderivatives (FTOC 1
        guarantees that it does!), but we can't write down any sort of
        nice formula for those antiderivatives.
      \end{solution}

    \end{enumerate}    
  \probonly{}

  \pspace{\newpage}

\item 
  \note{I will use this to try to motivate the method of substitution; we
    know that the second statement follows from the first by the Chain
    Rule, so it is reasonable to say $u = x^2$ and $du = 2x\,dx$ to
    ``match up'' the pieces of the two integrals.}

  Fill in the blank in each statement:

  \begin{enumerate}
  \item
    \begin{problem}
      $\displaystyle \int \cos(u)\, du = \fitb{1.25in}{\sin(u) + C}$
    \end{problem}

    \begin{solution}
      The derivative of $\sin(u)$ is $\cos u$.
    \end{solution}

  \item
    \begin{problem}
      $\displaystyle \int \cos(x^2) \cdot 2x\, dx =
      \fitb{1.25in}{\sin(x^2) + C}$ because $\dfrac{d}{dx} \sin(x^2) = 
      \fitb{1.25in}{\cos(x^2) \cdot 2x}$
    \end{problem}

    \begin{solution}
      By the Chain Rule, the derivative of $\sin(x^2)$ is $\cos(x^2) \cdot
      2x$.
    \end{solution}

  \end{enumerate}

\item 
  \note{In this first group of problems, letting $u$ be the inside of the
    composition works every time. I plan to do \#3a,b,c,d as a whole class.
    I'll have the six integrals at the top of the page in a box on the board
    and ask students ``Which of these is the underlying structure of the
    integral we're facing?'' I'll use colored chalk to write\\
    \null\quad Let $u=$\\
    \null\quad Then $du=$}

  Evaluate the following integrals using substitution.


  \begin{enumerate}

  \item
    \begin{problem}[\vfill]
      $\displaystyle \int e^{\sin x} \cos x\,dx$
    \end{problem}

    \begin{solution}
      It looks like a good substitution is \sub{$u = \sin x$}, since $\sin
      x$ is the inner function in the composition $e^{\sin x}$, and we
      also see \subd{$du = \cos x\,dx$} in the integral.  Then,
      \begin{align*}
        \int e^{\sub{\sin x}} \subd{\cos x\,dx} &= \int e^\sub{u}\,\subd{du} \\
        &= e^u + C \\
        &= e^{\sin x} + C
      \end{align*}
      \checkmark\
        Remember that you can always check your answer by differentiating
        it! 
    \end{solution}

  \item
    \begin{problem}[\vfill]
      $\displaystyle \int \sin x \cos^5 x\,dx$
    \end{problem}

    \begin{solution}
      It looks like a good substitution is \sub{$u = \cos x$} because it's
      the inner function in the composition $\cos^5 x$, and we see something related to $du=-\sin x\,dx$ in the integral, namely
      \subd{$-du = \sin x\,dx$}.
      \begin{align*}
        \int \sin x \cos^5 x\,dx & = \int (\subd{\sin x}) (\sub{\cos
          x})^5\,\subd{dx} \\ 
        &= \int \sub{u}^5 (\subd{-du}) \\
        &= - \int u^5\,du \\
        & = - \frac{1}{6} u^6 + C \\
        & = \boxed{- \frac{1}{6} \cos^6 x + C}
      \end{align*}
    \end{solution}

  \item

    \note{This is the first definite integral, and we'll talk about 2 ways
      to deal with the endpoints: one is to ignore them, find the
      indefinite integral as usual, and then plug the endpoints into that.
      The other is to change the endpoints to be in terms of $u$.}

    \begin{problem}[\nstretch{2}]
      $\displaystyle \int_0^{\pi/2} \sin x \cos^5 x \, dx$
    \end{problem}

    \begin{solution}
    There are two ways to approach this. One is to use
    the general antiderivative of $\sin x \cos^5 x$ that we just found,
      $-\frac{1}{6}\cos^6 x +C$. With this, we can find
      \begin{align*}
          \int_0^{\pi/2} \sin x \cos^5 x \,dx &= 
          \left. -\frac{1}{6}\cos^6 x\right|_0^{\pi/2} \\
          &= -\frac{1}{6}\cos^6\left(\frac{\pi}{2}\right) - 
          \left(-\frac{1}{6}\cos^6(0)\right) \\
          &= \boxed{\frac{1}{6}}
      \end{align*}
      Another way is to use the same substitution as before, \sub{$u = \cos x$}
      and \subd{$du = - \sin x\,dx$}, but to also rewrite the bounds in
      terms of $u$. When $x=0$, $u=\cos(0)=1$. When $x=\pi/2$,
      $u=\cos(\pi/2)=0$. 
      \begin{align*}
        \int_0^{\pi/2} \sin x \cos^5 x\,dx 
        & = \int_{x=0}^{x=\pi/2} (\subd{-\sin x}) (\sub{\cos
          x})^5\,\subd{dx} \\ 
        &= \int_{u=1}^{u=0} \sub{u}^5 (\subd{-du}) \\
        &= - \int_1^0 u^5\,du \\
        &= \int_0^1 u^5\,du\\
        &= \left.\frac{1}{6} u^6\right|_0^1 \\
        &= \boxed{\frac{1}{6}}
      \end{align*}
    \end{solution}

  \item 
    \begin{problem}[\vfill]
      $\displaystyle \int_1^2 \frac{x}{x^2+5}\,dx$
    \end{problem}

    \begin{solution}
        It looks like a good substitution is \sub{$u = x^2+5$} because it's
      the inner function in the composition $\dfrac{1}{x^2+5}$, and we can
      get \subd{$du = 2x\,dx$} in the integral by multiplying by 
      $\frac{1}{2}\cdot2$. The new bounds will be from $u=6$ to $u=9$.
      \begin{align*}
          \int_1^2 \frac{x}{x^2+5}\,dx &= 
          \frac{1}{2}\int_{x=1}^{x=2} 
          \frac{\subd{2x}}{\sub{x^2+5}}\,\subd{dx} \\
          &= \frac{1}{2}\int_{u=6}^{u=9} \frac{1}{\sub{u}}\,\subd{du} \\
          &= \frac{1}{2} \left[\ln|u|\right]_6^9 \\
          &= \boxed{\frac{1}{2}(\ln 9 - \ln 6)}\\
      \end{align*}
    \end{solution}

  \item
    

    \begin{problem}[\vfill]
      $\displaystyle \int \frac{\sin (\ln x)}{5x}\,dx$
    \end{problem}

    \begin{solution} % originally by Henry Shin, modified by Janet
      It looks like a good substitution is \sub{$u = \ln x$}, since it's
      the inner function in the composition $\sin(\ln x)$, and we also see
      \subd{$du = \frac{1}{x}\,dx$} in the integral.  So,
      \begin{align*}
        \int \frac{\sin (\sub{\ln x})}{5\subd{x}}\,\subd{dx} 
        & =  \int \frac{\sin \sub{u}}{5} \,\subd{du} \\
        & =  -\frac{1}{5} \cos u + C \\
        &= \boxed{-\frac{1}{5} \cos(\ln x) +C}
      \end{align*}
    \end{solution}

  \item
    \note{Occasionally students want to rewrite this as $\displaystyle
      \int_1^4 \frac{1}{3} \sqrt{u}\,du$; they know in their head that $x$
      goes from $1$ to $4$, and they get the right answer at the end, so
      they don't really see why we object to them writing $\displaystyle
      x^2 \sqrt{1 + x^3}\,dx = \int_1^4 \frac{1}{3} \sqrt{u}\,du$.  Have
      them look back at {{{{\#1}}(a)}} to see that
      $\displaystyle \int_1^4 \frac{1}{3} \sqrt{u}\,du$ has its own
      meaning, and that although what they write may make sense to them,
      it wouldn't make sense to anyone else.}

    \begin{problem}[\vfill]
      $\displaystyle \int_1^4 x^2 \sqrt{1 + x^3}\,dx$
    \end{problem}

    \begin{solution}
      It looks like a good substitution is \sub{$u = 1 + x^3$} since it's
      the inner function in the composition $\sqrt{1 + x^3}$, and we also
      see something related to $du = 3x^2\,dx$, namely \subd{$\frac{1}{3}
        du = x^2\,dx$}.  When $x = 1$, $u = 1 + 1^3 = 2$; when $x = 4$, $u
      = 1 + 4^3 = 65$.  So,
      \begin{align*}
        \int_1^4 \subd{x^2} \sqrt{\sub{1 + x^3}}\,\subd{dx} &=
        \int_2^{65} \sqrt{\sub{u}} \cdot \subd{\frac{1}{3}\, du} \\
        &= \frac{1}{3} \int_2^{65} u^{1/2}\,du \\
        & = \left. \frac{2}{9} u^{3/2} \right|_2^{65} \\
        & = \boxed{\frac{2}{9} \left( 65^{3/2} - 2^{3/2} \right)}
      \end{align*}
      \emph{Note on notation:} If you prefer not to change the bounds
      of the integral to be in terms of $u$, your first line should
      instead read      
      \[ \int_1^4 {x^2} \sqrt{{1 + x^3}}\,{dx} =
      \int_{{\color{red}x=\color{black}1}}^{\color{red}x=\color{black}4}
      \sqrt{u} \cdot {\frac{1}{3}\, du}.\] It is \color{red}\textbf{not}
      \color{black} correct to write \[ \int_1^4 {x^2} \sqrt{{1 +
          x^3}}\,{dx} = \int_1^4 \sqrt{u} \cdot {\frac{1}{3}\, du}.\]
      After all, we saw in {{{{\#1}}(a)}} that $\displaystyle
      \int_1^4 \sqrt{u} \cdot {\frac{1}{3}\, du}$ is equal to
      $\dfrac{14}{9}$, and we've just seen that $\displaystyle \int_1^4
      {x^2} \sqrt{{1 + x^3}}\,{dx} \neq \frac{14}{9}$.
    \end{solution}

  \end{enumerate}

  \pspace{\newpage}

\item 
  \emph{Examples where substitution is useful even though there isn't an
    obvious composition.}

  \begin{enumerate}
  \item

    \note{This exact (classic) problem reappears on their homework. Don't
      mention this to them! We want them to learn to recognize it on their
      own.}

    \begin{problem}[\vfill]
      $\displaystyle \int \tan x\,dx$
    \end{problem}

    \begin{solution}
      Let's first rewrite $\tan x$ as $\displaystyle \frac{\sin x}{\cos
        x}$, so that we are trying to evaluate $\displaystyle \int
      \frac{\sin x}{\cos x}\,dx$.  It looks like \sub{$u = \cos x$} is a
      promising substitution because we also see something related to
      $du = - \sin x\,dx$, namely \subd{$-du = \sin x\,dx$}:
      \begin{align*}
        \int \frac{\subd{\sin x}}{\sub{\cos x}}\,\subd{dx} &= \int
        \frac{1}{\sub{u}} \left( \subd{-du} \right) \\
        &= - \int \frac{1}{u}\,du \\
        &= - \ln |u| + C \\
        &= \boxed{- \ln |\cos x| + C}
      \end{align*}
    \end{solution}

  \item

    \begin{problem}[\vfill]
      $\displaystyle \int_{1/e}^{e^3} \frac{4 \ln x}{x}\,dx$
    \end{problem}

    \begin{solution}
      It looks like \sub{$u = \ln x$} is a promising substitution because
      we also see \subd{$du = \frac{1}{x}\,dx$} in the integral:
      \begin{align*}
        \int_{1/e}^{e^3} \frac{4 \sub{\,\ln x}}{\subd{x}}\,\subd{dx} &=
        \int_{-1}^3 
        4\sub{u}\,\subd{du} \\ 
        &= 2u^2 \Big|_{-1}^3 \\
        &= 2 \cdot 3^2 - 2 \cdot (-1)^2 \\
        &= \boxed{16}
      \end{align*}
    \end{solution}

  \end{enumerate}

\item 
  \begin{problem}[\vfill]
    One of the following integrals can easily be evaluated using
    substitution.  Which one?  Evaluate it.
    
      \begin{enumerate}
      \item $\displaystyle \int e^{-x^2}\,dx$
      \item $\displaystyle \int xe^{-x^2}\,dx$
      \end{enumerate}
    
  \end{problem}

  \begin{solution}
    Both integrals involve a composition with \sub{$u = -x^2$} as the
    inner function, but only the second integral has something related to
    $du = - 2x\,dx$, namely \subd{$-\frac{du}{2} = x\,dx$}.  That's the
    one we can easily evaluate with substitution: 
    \begin{align*}
      \int \subd{x}e^{\sub{-x^2}}\,dx &= \int e^\sub{u} \left( \subd{-
          \frac{du}{2}} \right) \\
      &= - \frac{1}{2} \int e^u\,du \\
      &= - \frac{1}{2} e^u + C \\
      &= \boxed{- \frac{1}{2} e^{-x^2} + C}
    \end{align*}
    The other indefinite integral, $\displaystyle \int e^{-x^2}\,dx$,
    cannot be easily evaluated using substitution.  In fact, like in
    {{{{\#1}}(c)}}, the antiderivatives of
    $e^{-x^2}$ cannot be expressed as a finite sum of familiar functions
    (like exponentials, logs, trig functions, etc.), so $\displaystyle
    \int e^{-x^2}\,dx$ cannot  be easily evaluated with any method!
  \end{solution}

% \item 
%   \note{I don't expect to get to this problem; it's just here as a
%     challenge for any student who can race through the rest of the
%     worksheet to get here.} 

%   Suppose $g$ is a mysterious function, and the only thing you
%   know about it is that $\displaystyle \int_0^{12} g(x)\,dx = \pi$.  For
%   each of the following integrals, decide whether you have enough
%   information to evaluate the integral; if so, evaluate it. 
%   \probonly{}
%     \begin{enumerate}
%     \item
%       \begin{problem}
%         $\displaystyle \int_0^{12} g(3x)\,dx$
%       \end{problem}

%       \begin{solution}
%         There is \fbox{not enough information}.  (As you can see in the
%         next part, substituting $u = 3x$ doesn't help.)  We can understand
%         why graphically; suppose $g$ looks like this.  The only
%         information we're given is that the green area is $\pi$:
%         \begin{center}
%           \begin{tikzpicture}[declare
%             function={f(\x)=(\x<=4*pi)*(\x/1.2^\x+1.5)+
%               (\x>4*pi)*(0.5*(5.542254215512511 + abs(sin(deg(0.5*\x)))));}]
%             \begin{axis}[ymin=0, xtick={12}, ytick=\empty, width=2.5in,
%               height=1.5in, clip=false] 
%               \addplot+[domain=0:12, plusarea] { f(x) } \closedcycle; %
%               \addplot+[domain=0:12*pi, samples=121] { f(x) }
%               node[right] {$y = g(x)$}; %
%             \end{axis}
%           \end{tikzpicture}
%         \end{center}
%         The graph of $g(3x)$ is stretched horizontally to be $\frac{1}{3}$
%         as wide, so here's how we picture $\displaystyle \int_0^{12}
%         g(3x)$:
%         \begin{center}
%           \begin{tikzpicture}[declare
%             function={f(\x)=(\x<=4*pi)*(\x/1.2^\x+1.5)+
%               (\x>4*pi)*(0.5*(5.542254215512511 + abs(sin(deg(0.5*\x)))));}]
%             \begin{axis}[ymin=0, xtick={12}, ytick=\empty, width=2.5in,
%               height=1.5in, clip=false] 
%               \addplot+[domain=0:12, fillregion] { f(3*x) } \closedcycle; %
%               \addplot+[domain=0:12*pi, samples=121] { f(3*x) }
%               node[right] {$y = g(3x)$}; %
%             \end{axis}
%           \end{tikzpicture}
%         \end{center}
%         But there's no way for us to find this area because we can't
%         express it simply in terms of the green area above, which is the
%         only area we know. 
%       \end{solution}

%     \item
%       \begin{problem}
%         $\displaystyle \int_0^4 g(3x)\,dx$
%       \end{problem}

%       \begin{solution}
%         We can evaluate this by substituting \sub{$u = 3x$}; then, $du =
%         3\,dx$, so \subd{$\frac{du}{3} = dx$}.  Since $x$ goes from $0$ to
%         $4$, $u = 3x$ goes from $0$ to $12$; thus,
%         \begin{align*}
%           \int_0^{4} g(\sub{3x})\,\subd{dx} &= \int_0^{12} g
%           (\sub{u})\,\subd{\frac{du}{3}} \\
%           &= \frac{1}{3} \int_0^{12} g(u)\,du \\
%           &= \boxed{\frac{\pi}{3}}
%         \end{align*}
%         This should also make sense graphically; the graph of $g(3x)$ is
%         obtained by stretching the graph of $g(x)$ horizontally by a
%         factor of $\frac{1}{3}$:
%         \begin{center}
%           \begin{tikzpicture}[declare
%             function={f(\x)=(\x<=4*pi)*(\x/1.2^\x+1.5)+
%               (\x>4*pi)*(0.5*(5.542254215512511 + abs(sin(deg(0.5*\x)))));}]
%             \begin{axis}[ymin=0, xtick={12}, ytick=\empty, width=2.5in,
%               height=1.5in, clip=false] 
%               \addplot+[domain=0:12, plusarea] { f(x) } \closedcycle; %
%               \addplot+[domain=0:12*pi, samples=121] { f(x) }
%               node[right] {$y = g(x)$}; %
%             \end{axis}
%           \end{tikzpicture}
%           \hspace{0.5in}
%           \begin{tikzpicture}[declare
%             function={f(\x)=(\x<=4*pi)*(\x/1.2^\x+1.5)+
%               (\x>4*pi)*(0.5*(5.542254215512511 + abs(sin(deg(0.5*\x)))));}]
%             \begin{axis}[ymin=0, xtick={4}, ytick=\empty, width=2.5in,
%               height=1.5in, clip=false] 
%               \addplot+[domain=0:4, fillregion] { f(3*x) } \closedcycle; %
%               \addplot+[domain=0:12*pi, samples=121] { f(3*x) }
%               node[right] {$y = g(3x)$}; %
%             \end{axis}
%           \end{tikzpicture}
%         \end{center}
%         Therefore, the area under $g(3x)$ from $x = 0$ to $x = 4$ should
%         be $\frac{1}{3}$ as large as the area under $g(x)$ from $x = 0$ to
%         $x = 12$. 
%       \end{solution}

%     \item
%       \begin{problem}
%         $\displaystyle \int_0^{144} g(x^2)\,dx$
%       \end{problem}

%       \begin{solution}
%         There is \fbox{not enough information}.  (If you try substituting
%         $u = x^2$, you instead end up with the integral in
%         {{{{\#6}}(d)}}.) 
%       \end{solution}

%     \item
%       \begin{problem}
%         $\displaystyle \int_0^{2 \sqrt{3}} x g(x^2)\,dx$
%       \end{problem}

%       \begin{solution}
%         We can evaluate this using the substitution \sub{$u = x^2$}; then,
%         $du = 2x\,dx$, so \subd{$\frac{du}{2} = x\,dx$}.  As $x$ goes from
%         $0$ to $2 \sqrt{3}$, $u = x^2$ goes from $0^2 = 0$ to $\left(2
%         \sqrt{3} \right)^2 = 12$:
%         \begin{align*}
%           \int_0^{2 \sqrt{3}} \subd{x} g(\sub{x^2})\,\subd{dx} &=
%           \int_0^{12} g(\sub{u}) \subd{\frac{du}{2}} \\
%           &= \frac{1}{2} \int_0^{12}
%           g(u)\,du \\
%           &= \boxed{\frac{\pi}{2}}
%         \end{align*}
%       \end{solution}

%     \item
%       \begin{problem}
%         $\displaystyle \int_0^{144} x g(x^2)\,dx$
%       \end{problem}      

%       \begin{solution}
%         There is \fbox{not enough information}.
%       \end{solution}

%     \end{enumerate}
%     \probonly{}

%   \pspace{\vfill}

\item  \note{Just additional practice}  

  Evaluate the following integrals. 
  \begin{enumerate}

  \item
    \begin{problem}[\vfill]
      $\displaystyle \int_0^1 \sqrt{5x + 4}\,dx$
    \end{problem}

    \begin{solution}
      It looks like \sub{$u = 5x + 4$} is a good substitution since it's
      the inner function in the composition $\sqrt{5x + 4}$.  Then, $du =
      5\,dx$, so \subd{$\frac{du}{5} = dx$}.  As $x$ goes from $0$ to $1$,
      $u = 5x + 4$ goes from $u = 5 \cdot 0 + 4 = 4$ to $5 \cdot 1 + 4 =
      9$:
      \begin{align*}
        \int_0^1 \sqrt{\sub{5x + 4}}\,\subd{dx} &= \int_4^9
        \sqrt{\sub{u}}\,\subd{\frac{du}{5}} \\
        &= \frac{1}{5} \int_4^9 u^{1/2}\,du \\
        &= \left. \frac{1}{5} \cdot \frac{2}{3} u^{3/2} \right|_4^9 \\
        &= \frac{2}{15} \left( 9^{3/2} - 4^{3/2} \right) \\
        &= \frac{2}{15} (27 - 8) \\
        &= \boxed{\frac{38}{15}}
      \end{align*}
    \end{solution}

  \item
    \begin{problem}[\vfill]
      $\displaystyle \int \frac{e^x}{\sqrt{2e^x + 5}}\,dx$
    \end{problem}

    \begin{solution} % originally by Henry Shin, modified by Janet
      It looks like \sub{$u = 2e^x + 5$} is a good substitution since it's
      the inner function in the composition $\sqrt{2e^x + 5}$, and we also
      see something related to $du = 2 e^x\,dx$ in the integral, namely
      \subd{$\frac{du}{2} = e^x\,dx$}: 
      \begin{align*}
        \int \frac{\subd{e^x}}{\sqrt{\sub{2e^x + 5}}}\,\subd{dx} & =
        \int \frac{1}{\sqrt{\sub{u}}} \cdot \subd{\frac{du}{2}} \\
        & = \int \frac{1}{2} u^{-1/2}\,du \\
        & =  u^{1/2} + C \\
        & = \boxed{\sqrt{2e^x + 5} + C}
      \end{align*}
    \end{solution}

  % \item
  %   \begin{problem}[\stretch{2}]
  %     $\displaystyle \int \frac{x+1}{\sqrt[3]{3x^2 + 6x + 5}}\,dx$        
  %   \end{problem}

  %   \begin{solution}
  %     It may not be entirely clear what to substitute here, but a
  %     reasonable first try is \sub{$u = 3x^2 + 6x + 5$}, since it's the
  %     inner function in the composition $\sqrt[3]{3x^2 + 6x + 5}$.  With
  %     this choice, $du = (6x + 6)\,dx = 6(x + 1)\,dx$, which looks
  %     promising since we see something related to this, namely
  %     \subd{$(x+1)\,dx = \frac{1}{6}\,du$}.  Then,
  %     \begin{align*}
  %       \int \frac{\subd{x + 1}}{\sqrt[3]{\sub{3x^2 + 6x +
  %             5}}}\,\subd{dx} &= \int 
  %       \frac{1}{\sqrt[3]{\sub{u}}}\,\subd{\frac{du}{6}} \\
  %       &= \frac{1}{6} \int u^{-1/3}\,du \\
  %       &= \frac{1}{6} \cdot \frac{3}{2} u^{2/3} + C \\
  %       &= \frac{1}{4} u^{2/3} + C \\
  %       &= \boxed{\frac{1}{4} (3x^2 + 6x + 5)^{2/3} + C}
  %     \end{align*}
  %   \end{solution}

  \end{enumerate}

\end{enumerate}

\end{document}